\documentclass[11pt,a4paper]{article}

\usepackage[utf8]{inputenc}
\usepackage[T1]{fontenc}
\usepackage[french]{babel}
\usepackage{amsmath}
\usepackage{amsfonts}
\usepackage{graphicx}
\usepackage{booktabs}
\usepackage{geometry}
\usepackage{hyperref}

\geometry{margin=2.5cm}

\title{Hybridation de Métaheuristiques pour le Problème du Voyageur de Commerce (TSP)}
\author{Master Intelligence Artificielle et Recherche Opérationnelle}

\begin{document}

\maketitle

\begin{abstract}

Le problème du voyageur de commerce (Travelling Salesman Problem, TSP) constitue
l’un des problèmes d’optimisation combinatoire les plus étudiés dans la littérature.
Sa résolution exacte devient rapidement intractable lorsque le nombre de villes
augmente. Les métaheuristiques représentent alors une alternative efficace pour
obtenir des solutions de haute qualité dans un temps de calcul raisonnable.

Dans ce projet, nous étudions plusieurs méthodes de résolution pour le TSP :
une heuristique gloutonne, l’algorithme Threshold Accepting et le recuit simulé.
Nous proposons également une approche hybride combinant ces méthodes afin
d’améliorer les performances globales. Les expériences sont réalisées sur
plusieurs instances de test et les résultats obtenus sont comparés en termes
de qualité de solution et de temps de calcul.

\end{abstract}

\section{Introduction}

Le problème du voyageur de commerce (TSP) consiste à déterminer un cycle de coût
minimal passant exactement une fois par chaque ville d’un ensemble donné.
Ce problème apparaît dans de nombreuses applications réelles telles que :

\begin{itemize}
\item la planification de tournées logistiques,
\item l’optimisation des réseaux de transport,
\item la robotique et la planification de trajectoires,
\item l’optimisation des circuits électroniques.
\end{itemize}

Le TSP appartient à la classe des problèmes NP-difficiles.
Ainsi, les méthodes exactes deviennent rapidement impraticables
lorsque la taille du problème augmente.

Pour cette raison, les métaheuristiques sont largement utilisées
afin d’obtenir des solutions de bonne qualité en un temps raisonnable.
Dans ce travail, nous étudions plusieurs approches heuristiques et
métaheuristiques pour résoudre ce problème.

\section{Formulation du problème}

Soit un ensemble de villes :

\[
V = \{1,2,\dots,n\}
\]

et une matrice de distances :

\[
D = (d_{ij})
\]

où $d_{ij}$ représente la distance entre la ville $i$ et la ville $j$.

Le problème du voyageur de commerce consiste à déterminer un cycle
passant exactement une fois par chaque ville et minimisant la distance
totale parcourue.

Une solution du problème peut être représentée par une permutation
des villes :

\[
\pi = (\pi(1), \pi(2), \dots, \pi(n))
\]

où $\pi(i)$ désigne la $i$-ème ville visitée dans la tournée.

La distance totale associée à cette tournée est donnée par :

\[
f(\pi) = \sum_{i=1}^{n} d_{\pi(i)\pi(i+1)}
\]

avec la condition :

\[
\pi(n+1) = \pi(1)
\]

qui garantit le retour à la ville de départ.

L’objectif est donc de déterminer la permutation $\pi$ minimisant
la distance totale parcourue :

\[
\min_{\pi} f(\pi)
\]

Ce problème appartient à la classe des problèmes
d’optimisation combinatoire NP-difficiles, ce qui
justifie l’utilisation de méthodes heuristiques et
métaheuristiques pour obtenir des solutions de
bonne qualité dans un temps de calcul raisonnable.

\section{Méthodes utilisées}

\subsection{Heuristique du plus proche voisin}

L’heuristique du plus proche voisin est une méthode gloutonne
qui construit une solution progressivement.

L’algorithme fonctionne de la manière suivante :

\begin{enumerate}
\item choisir une ville de départ,
\item sélectionner la ville non visitée la plus proche,
\item répéter jusqu’à ce que toutes les villes soient visitées.
\end{enumerate}

Cette méthode est rapide mais peut produire des solutions de qualité
variable.

\subsection{Threshold Accepting}

Threshold Accepting est une métaheuristique introduite comme
alternative déterministe au recuit simulé.

L’idée consiste à accepter certaines solutions moins bonnes,
à condition que la dégradation ne dépasse pas un seuil donné.

Soit :

\[
\Delta = f(x') - f(x)
\]

où $x'$ est un voisin de la solution courante $x$.

La solution est acceptée si

\[
\Delta \leq S
\]

où $S$ représente un seuil décroissant au cours de l’algorithme.

Cette stratégie permet d’éviter les minima locaux.

\subsection{Recuit simulé}

Le recuit simulé est inspiré du processus physique de refroidissement
des métaux.

La probabilité d’accepter une solution moins bonne est donnée par :

\[
P = e^{-\Delta/T}
\]

où :

\begin{itemize}
\item $T$ est la température,
\item $\Delta$ est la variation de la fonction objectif.
\end{itemize}

Lorsque la température diminue progressivement,
l’algorithme devient de plus en plus sélectif.

\section{Hybridation proposée}

Afin d’améliorer les performances globales,
nous proposons une approche hybride combinant plusieurs méthodes.

La stratégie adoptée est la suivante :

\begin{enumerate}
\item génération d’une solution initiale par l’heuristique gloutonne,
\item amélioration locale par l’algorithme Threshold Accepting,
\item optimisation finale par recuit simulé.
\end{enumerate}

Cette hybridation permet de combiner :

\begin{itemize}
\item la rapidité de l’heuristique gloutonne,
\item la capacité d’exploration locale de Threshold Accepting,
\item la diversification du recuit simulé.
\end{itemize}

\section{Expérimentations}

Les expériences ont été réalisées sur plusieurs instances
du problème du voyageur de commerce.

Les algorithmes ont été implémentés en Python
et exécutés sur différentes tailles d’instances.

Les critères de comparaison sont :

\begin{itemize}
\item la distance totale du tour,
\item le temps d’exécution.
\end{itemize}

\begin{center}
\begin{tabular}{lccc}
\toprule
Méthode & Distance & Temps (s) \\
\midrule
Greedy & 2760 & 0.01 \\
Threshold Accepting & 2410 & 0.35 \\
Simulated Annealing & 2300 & 0.60 \\
Hybrid & \textbf{2200} & 0.70 \\
\bottomrule
\end{tabular}
\end{center}

On observe que l’approche hybride permet d’obtenir
les meilleures solutions.

\section{Discussion}

Les résultats montrent que :

\begin{itemize}
\item les heuristiques simples sont rapides mais peu précises,
\item les métaheuristiques améliorent significativement la solution,
\item l’hybridation permet de combiner les avantages de plusieurs méthodes.
\end{itemize}

Cette approche est particulièrement intéressante pour les problèmes
d’optimisation combinatoire de grande taille.

\section{Conclusion}

Dans ce travail, nous avons étudié plusieurs méthodes
pour résoudre le problème du voyageur de commerce.

Nous avons implémenté et comparé :

\begin{itemize}
\item une heuristique gloutonne,
\item l’algorithme Threshold Accepting,
\item le recuit simulé,
\item une approche hybride.
\end{itemize}

Les expériences réalisées montrent que l’hybridation
permet d’améliorer la qualité des solutions.

Des perspectives d’amélioration incluent :

\begin{itemize}
\item l’utilisation d’opérateurs de voisinage plus avancés,
\item l’intégration d’algorithmes génétiques,
\item l’utilisation de méthodes d’apprentissage automatique.
\end{itemize}

\end{document}
